% results.tex

We evaluated two forms of training-free conditioning across six protein families
(Table~\ref{tab:cross-family}). Kunitz domains served as the primary test case, and four
additional Pfam seed families---SH3, WW, Homeobox, and Forkhead---tested whether the findings
generalized across family size and PCA geometry. The $\omega$-conotoxin O-superfamily provided
a distinct peptide example associated with the Cav2.2 calcium
channel~\cite{omegaConotoxinSAR2006}. In every family, we
compared hard curation, which rebuilds the memory from a designated subset, with multiplicity
weighting, which retains the full-family memory but increases the designated patterns' weight.
No activity measurements were available for these families, so every designated subset was
defined from sequence alone: the sequences carrying a chosen residue class at one alignment
column for the Pfam families, and a curated list of accessions for the $\omega$-conotoxins.
Each designation was therefore a sequence feature and not a measured function.

\DisplayCrossFamily

We first asked whether hard curation could transfer a selected marker while preserving the
broader family pattern (Fig.~\ref{fig:scaling}). Kunitz domains (PF00014; $K=99$, $L=53$ aligned positions) provided an
interpretable test: their P1 residue influences which protease they inhibit, and Lys or Arg at
P1 is associated with trypsin specificity~\cite{laskowskiProteinInhibitors1980}. We designated the
32 K/R-at-P1 sequences as marker-positive and the remaining 67 as marker-negative.
Separate memories built from these two subsets produced the sharp result expected when the
input was selected on the marker itself. All 930 sequences generated from the marker-positive
memory carried K/R at P1, and none generated from the marker-negative memory did; full-family
generation produced 38\% K/R compared with 32\% in the input family. The conditioning remained
localized rather than erasing the family pattern. The top marker-positive sequences retained
all 11 highly conserved positions and all six cysteines
while carrying 19--26 substitutions concentrated at variable sites. Per-position Shannon
entropy correlated with that of the stored multiple sequence alignment (MSA) for both full-family
generation ($r=0.988$) and marker-positive conditioning ($r=0.932$), with the largest change
at P1 (Fig.~\ref{fig:sequence-analysis}). The result also held when the designated input was
subsampled from 32 sequences down to three: P1 retention remained 100\% at every tested size,
whereas pairwise diversity rose and Kullback--Leibler (KL) divergence fell as more inputs were
added (Fig.~\ref{fig:scaling}). Thus, even three inputs were sufficient to fix the selected
marker, while additional inputs coincided with a broader PCA subspace and better distributional fit.

\DisplayHardCuration

While hard curation provides a clear endpoint, it discards the background sequences and changes
the PCA basis. We therefore asked whether multiplicity weighting could tune Kunitz generation
continuously on the full-family basis (Fig.~\ref{fig:phase-transition}). The canonical
replicated analysis swept the multiplicity ratio $\rho$ between designated and background
weights from 1 (no bias) to 500 (strong bias); at each value, we recomputed the entropy-crossover temperature
$\beta^{*}$ and generated five independent libraries of 620 sequences (Table~\ref{tab:per-family-rho}).
We then compared the attention the sampler placed on designated patterns with the fraction of
decoded sequences carrying the P1 marker. Raising $\rho$ moved the attention, but the decoded
fraction followed only partway. At $\rho=500$, the weights called for 99.6\% of the sampler's
equilibrium mixture to sit on the designated Kunitz patterns. Every designated pattern carries
K/R at P1 and no background pattern does, so a library that carried the weighting through
unchanged would show K/R at P1 in 99.6\% of its sequences. The decoded fraction reached only
\KunitzFobsFiveHundred{} $\pm$ \KunitzFobsFiveHundredSD{}. This discrepancy between the target
marker fraction and the decoded one is the calibration gap $\Delta$
(Eq.~\ref{eq:decoder-pushforward}).
Increasing $\rho$ also produced modest global changes:
sequence diversity decreased from 0.56 to 0.51, KL divergence increased from 0.007 to 0.017,
and $\beta^{*}$ shifted from 4.4 at $\rho=1$ to 9.3 at $\rho=1{,}000$. Over that range the
effective pattern count, a measure of how concentrated the multiplicity
weights have become, fell from 99 to 32.1, approaching the concentration expected when weight
is confined to the 32 designated patterns. The decoded marker fraction nevertheless stayed far
below its target.

To locate the shortfall, we compared the analytic designated mixture share, empirical attention,
and decoded marker frequency. Across 17 values of $\rho$ from 1 to 1{,}000, mean designated
attention remained within 1.2 percentage points of the analytic share. The marker shortfall
therefore appeared downstream of latent weighting, where finite-temperature variation,
unit-normalized PCA reconstruction, and argmax decoding act together. Designated and background Kunitz sequences overlap in the
80-dimensional PCA basis (separation index $S=\KunitzSeparation$), so a decoded sequence need
not identify the component that generated its latent state. An exact-equilibrium comparison also
indicated that finite-step sampling was not the main source of the shortfall. At $\rho = 1$, 10,
and 500, we compared 1{,}640 independent draws from the
closed-form Gaussian mixture model with 1{,}640 retained unadjusted Langevin (ULA) states under
the same $\beta^{*}$ and decoder. At $\rho=500$, exact and ULA sampling produced nearly
identical decoded P1~K/R fractions of 0.604 and 0.597, and their decoded amino acid compositions differed by a KL
divergence of only $1.19\times10^{-4}$. Removing finite-step sampling error therefore moved the
decoded fraction by 0.007, against a shortfall of 0.39 from the 0.996 target. At all three values
of $\rho$, the fraction of states whose most probable mixture component was a designated
pattern differed by at most 0.026 between the two samplers, and the decoded marker fractions by
at most 0.016. The calibration gap therefore persisted at exact equilibrium and arose mainly
after latent sampling, in the combined reconstruction and decoding path.

\DisplayMultiplicity

The Kunitz result suggested that conditioning should transfer most effectively when the PCA
representation separates the designated and background subsets. We tested this hypothesis on
four additional Pfam families, tracking Trp at the peptide-groove column for SH3 (PF00018;
$K=55$), a specificity-loop residue for WW (PF00397; $K=420$), Gln at position~50 of the
recognition helix for Homeobox (PF00046; $K=136$), and His or Asn at the H3 recognition-helix
position for Forkhead (PF00250; $K=246$). Together, the five Pfam splits sampled distinct
PCA geometries at comparable conditioning strengths (Table~\ref{tab:per-family-rho}). The
separation index, larger when PCA resolves the two subsets more cleanly, ranged from
\WWSeparation{} for WW to \HomeoboxSeparation{} for Homeobox.
The attention tracked its target across the panel: over all six replicated family
sweeps, the largest absolute deviation between the attention on designated patterns and the
analytic designated share was \AttnMaxDevPts{} percentage points (\AttnMaxDevFamily{},
$\rho=\AttnMaxDevRho$). What varied was the decoded marker response.
Calibration gaps decreased with separation
overall, although not monotonically, and an exploratory five-family fit gave
$\Delta\approx\RelationIntercept{} \RelationSlope\,S$ with $R^{2}=\RelationRsq$
(Fig.~\ref{fig:separation-vs-gap}). However, the small panel did not support a universal rule:
leave-one-family-out slopes ranged from \RelationLooSlopeMin{} to \RelationLooSlopeMax{}, and
$R^{2}$ ranged from \RelationLooRsqMin{} to \RelationLooRsqMax{}. We therefore treated the fit
as a preliminary association. The $\omega$-conotoxin family appears in
Fig.~\ref{fig:separation-vs-gap} but not in the fit. It is not a Pfam seed family, and its
accession-based designation is not identical to the Tyr13 readout. Its plotted difference is
therefore a marker response, not a calibration error in the same sense as the five Pfam points.
Hard curation retained the marker completely in the five Pfam families and at 97.9\% in the
$\omega$-conotoxin family.

\DisplaySeparation

The calibration gap raised a natural comparator question: would it also appear in a model that
generates residues without SA's PCA reconstruction and argmax step?
We fitted a multiplicity-weighted profile HMM to each cleaned alignment, weighting every
designated sequence by $\rho$ and every background sequence by 1, so the HMM saw the same
split and the same multiplicity ratios as SA (Table~\ref{tab:profile-hmm-benchmark}). At
$\rho=500$, the profile HMM recovered the tracked marker at frequencies of 0.988--1.000 in the
five Pfam families, essentially reaching the target marker fraction that SA fell short of. The
contrast is consistent with a bottleneck in SA's latent-to-sequence path, but it does not isolate
that mechanism because the two models generate sequences differently. A profile HMM estimates
match-state emissions directly, whereas SA samples in PCA space and then reconstructs and
decodes a sequence. The $\omega$-conotoxin HMM produced Tyr13 in 0.834 of
its sequences and was essentially tied with SA at 0.838. Here $f_{\mathrm{eff}}=0.996$ refers
to the accession-designated mixture share, not a Tyr13 target. Both methods landed near
the marker prevalence of the designated input itself: only 19 of the 23 designated conotoxin
sequences carry Tyr13, so a generator that reproduces the designated marginals would yield the
marker in $19/23 = 0.83$ of the output. That prevalence is a reference point rather than an
upper bound, because a decoder that sharpens toward the mode can exceed it, as hard curation
does below. Closer marker matching did not imply lower language-model scores. In the matched
Kunitz $\rho=500$ comparison, ESM2 assigned pseudo-perplexity
$\SARhoFiveHundredPPL\pm\SARhoFiveHundredPPLSD$ to SA and
$\HMMRhoFiveHundredPPL\pm\HMMRhoFiveHundredPPLSD$ to the weighted profile HMM, where a lower
value means the model finds a sequence more plausible. At the same multiplicity ratio, the HMM
came closer to the target marker fraction, while SA had lower ESM2 pseudo-perplexity.

\DisplayProfileHMM

We next asked whether the marker signal lost between latent weighting and sequence output could
be recovered by sampling more sharply, without changing the memory
(Fig.~\ref{fig:mask-betasweep}). Two separate controls act
here: $\rho$ sets the target marker fraction, whereas the inverse temperature $\beta$ sets how
sharply the sampler retrieves individual stored patterns rather than blending them. We report
$\beta$ as the ratio $\beta/\beta^{*}$, since the crossover $\beta^{*}$ at which blending gives
way to retrieval itself shifts with $\rho$, so the ratio measures sharpness relative to each
condition's own operating point. Increasing $\beta/\beta^{*}$ from 0.5 to 3.0 at
$\rho=10$, 50, and 200 raised the decoded Kunitz P1 marker fraction but reduced diversity. At
$\rho=200$, for example, P1~K/R rose from 0.56 at $\beta^{*}$ to 0.70 at $3\beta^{*}$ while
diversity fell from 0.52 to 0.46 (Supporting Information). The target at that $\rho$ was 0.99,
so tripling the sharpness recovered about a third of the shortfall and cost 0.06 in diversity.
That left two contributors to the remaining shortfall: attention on background patterns at
finite $\rho$, and the path from a latent component to a decoded sequence. We therefore gave background patterns
zero weight inside the softmax and swept $\beta$ (Table~\ref{tab:mask-recovery}). This hard
mask puts all attention on designated patterns while leaving finite-temperature latent spread,
PCA reconstruction, and argmax decoding unchanged. Any remaining shortfall is therefore
downstream of attention allocation.
Under the hard mask, P1~K/R was only 0.50 at its operating point, $\beta^{*}=3.23$,
and reached 1.00 only near $\beta/\beta^{*}\approx160$. Complete recovery therefore takes a
sampler running two orders of magnitude sharper than its own operating point. At that sharpness
the sampler approaches the stored memories and explores less broadly,
and novelty fell from 0.43 to 0.16 over the same range. \SIRef{app:matched-beta} reports a
matched-temperature comparison of the mask against multiplicity weighting that reaches the same
conclusion. Together, these experiments showed that
sharper retrieval can compensate for latent-to-sequence loss, but only by sacrificing novelty.

\DisplayHardMask

Having established the conditioning mechanics on Pfam domains, we asked whether hard curation
could transfer features in a distinct disulfide-rich peptide family
(Table~\ref{tab:conotoxin-pharmacophore}). $\omega$-Conotoxins block
voltage-gated calcium channels, including Cav2.2~\cite{omegaConotoxinSAR2006}; the family
includes MVIIA (ziconotide), an approved non-opioid analgesic~\cite{ziconotideReview2004}.
Prior mutagenesis found Tyr13 essential for MVIIA and GVIA activity, making
Tyr13 a useful sequence readout here~\cite{kimTyr13Essential1995}. Lys2 and other basic residues
have also been studied as channel-binding determinants~\cite{satoBasicResidues1993}. We aligned 74 SwissProt
$\omega$-conotoxins and retained 26 positions after gap filtering. Residue numbers below denote
retained alignment columns mapped onto MVIIA, so the marker is retained column 13, corresponding
to MVIIA Tyr13. Of those 74 sequences, 23 accessions
appeared on the curated list and formed the designated subset. The remaining 51 formed the
background, for which we assumed no activity label.
Importantly, accession designation and Tyr13 were related but
distinct labels: they agreed for 64 of 74 sequences (86.5\%), with 19 Tyr13-positive sequences
in the designated set and six in the background. We therefore reported accession conditioning
and Tyr13 recovery separately; an accession-level audit identified which designations had
repository activity metadata and which were supported only by the input list.
Hard curation again shifted the selected sequence statistics. From 1{,}550 generated sequences
per condition, Tyr occurred at the marker position in 48.5\% of full-family generation versus
33.8\% of its input, and in 98.5\% of designated-subset generation versus 82.6\% of its input
(Table~\ref{tab:conotoxin-pharmacophore}). Designated-subset generation also contained 20.1\%
Lys or Arg, close to 19.7\% in its input and above the 12.0\% in full-family generation.
Multiplicity weighting transferred the marker in this family as well, and produced the largest
gain in the panel. Sweeping $\rho$ from 1 to 500 raised decoded Tyr13 from $0.463 \pm 0.040$ to
$0.838 \pm 0.035$ (Supporting Information), against 0.587 for Kunitz at the same ratio. That
endpoint sits near the 0.83 prevalence of Tyr13 among the 19 of 23 designated accessions that
carry it, which again marks where a generator reproducing the designated marginals would land
rather than a bound on what decoding can reach. The $\omega$-conotoxin split also has the
cleanest PCA separation in the panel
($S = \ConotoxinSeparation$, against \WWSeparation{}--\HomeoboxSeparation{} for the Pfam
families), qualitatively consistent with the preliminary association reported above; because
this family was held out of that fit, the agreement does not establish a predictive
relationship.

The broader sequence pattern was retained alongside these marker shifts. Frequency heatmaps
preserved Tyr13, the cysteine framework, and variable-loop diversity
(Fig.~\ref{fig:conotoxin-loop}); measured against the full-family input entropy profile,
per-position correlations were $r=0.776$ for designated-subset generation and $r=0.997$ for
full-family generation. Among the 50 modeled designated-subset sequences, the five
highest-pLDDT sequences retained all six highly conserved cysteine positions while introducing
3--9 substitutions at variable sites (Fig.~\ref{fig:conotoxin-sequence-analysis}).
Comparison with 12 published structure-activity relationship positions gave a consistent but
strictly sequence-level result. Designated-subset generation retained Tyr13 at 98.5\% and all
six invariant disulfide-forming cysteines. It also increased K/R at reported SAR
positions, including the Arg10-aligned column (83.4\% versus 69.6\% in the designated input) and
position~21 (52.4\% versus 47.8\%). These shifts describe residue frequencies; they do not test
mutation effects in the generated backgrounds~\cite{kimTyr13Essential1995,satoBasicResidues1993,lewisConotoxinSAR2012}.
The single highest-TM-score modeled sequence also aligned visually with the short experimental
reference backbone (Fig.~\ref{fig:fold-superposition}). Thus, without
retraining or a Cav2.2 target structure, hard curation transferred the designated accessions'
sequence statistics into a new library; it did not establish target binding.

\DisplayConotoxin

Finally, we asked whether marker control came at the expense of global sequence or structural
plausibility (\SIFigRef{fig:plddt-tmscore-scatter}). We scored 50 Kunitz sequences from each of
five sources: stored, full-family SA, marker-positive SA, marker-negative SA, and unconditional
profile-HMM emission. Across those sources, ESMFold assigned similar model scores to SA and
stored sequences (Table~\ref{tab:structure-validation}). Mean pLDDT was $90.4\pm1.2$ for
full-family SA, $90.7\pm1.1$ for marker-positive SA, $91.0\pm1.2$ for marker-negative SA, and
$89.3\pm3.2$ for stored sequences, with every SA sequence above 70. Corresponding TM-scores
against the experimental Kunitz structure (1BPI, chain~A) were $0.84\pm0.01$, $0.83\pm0.01$,
$0.85\pm0.01$, and $0.83\pm0.02$. Sequences emitted by the unconditional profile HMM, built with
default \texttt{hmmbuild} settings on the seed alignment rather than the weighted build used for
the matched benchmark, scored lower and more variably, with mean pLDDT $60.4\pm13.1$, 20\% above
70, and mean TM-score $0.60\pm0.19$. AlphaFold2-ptm provided a second structure predictor and
gave the same Kunitz pattern, with TM-scores of
0.83--0.85 at pLDDT 93--95 (Table~\ref{tab:structure-model-comparison})
~\cite{jumperHighlyAccurate2021,mirdita2022colabfold}. The $\omega$-conotoxins, at 24 to 35
residues, scored lower against their reference structure (1OMG, chain~A), but so did the stored
natural sequences scored alongside them. Under AlphaFold2-ptm, TM-scores were $0.35\pm0.11$ for stored
$\omega$-conotoxins, $0.35\pm0.13$ for full-family generation, and $0.47\pm0.02$ for
designated-subset generation, at pLDDT 68--80; ESMFold gave $0.43\pm0.07$, $0.47\pm0.03$, and
$0.47\pm0.03$ for the same three groups (Table~\ref{tab:structure-model-comparison}). The
designated-subset group had a higher mean TM-score than the stored group under both predictors.
The full-family group also had a higher mean under ESMFold; under AlphaFold2-ptm, its mean exceeded
the stored mean by only 0.0027, negligible relative to the within-group standard deviations.
Because every score was normalized to the 25-residue reference chain and the stored and generated
distributions overlapped, neither predictor showed a lower generated-group mean relative to the
stored comparison.
The low absolute values do not by themselves distinguish short-chain TM-score behavior from
structural differences among the peptides.

Sequence scoring with ESM2 also placed SA near the stored family distribution
(Table~\ref{tab:structure-validation}). Pseudo-perplexity was $4.78\pm0.64$ for full-family SA,
$4.72\pm0.60$ for marker-positive SA, and $4.97\pm0.98$ for stored sequences, compared with
$16.1\pm4.2$ for unconditional HMM emission. The weighted profile HMM of the matched benchmark,
fitted to the cleaned alignment without entropy weighting, instead scored
$\HMMRhoFiveHundredPPL\pm\HMMRhoFiveHundredPPLSD$, so the contrast drawn here is with
unconditional emission and not with profile HMMs in general.
Sequence-marker metrics from the same unconditional baseline confirmed the distinction among
generation, conditioning, and resampling. Across five HMMER3 emissions and five bootstrap
resamples~\cite{hmmer3}, full-family SA gave a P1~K/R fraction of $0.38\pm0.03$ and KL divergence
$0.0076\pm0.0009$, compared with $0.17\pm0.04$ and $0.027\pm0.002$ for HMM emission.
Bootstrap resampling matched composition most closely (KL $\approx0.001$) but produced no novel
sequences. Marker-positive SA, by contrast, reached P1~K/R $1.0\pm0.0$, a conditioning control
that neither unconditional baseline provided (Table~\ref{tab:structure-validation}). These
model-based evaluations did not demonstrate function. They indicate that marker control remained
compatible with family-level sequence plausibility and Kunitz-like predicted folds.

\DisplayModelValidation
